\subsection{Canonical and image restrictions in Albanese degree two}
\label{subsec:smallest-endpoint-strict-canonical}

The finite Stein factor and the Albanese image of the distinguished
curve impose related restrictions on the smallest endpoint.
The canonical bound is strict.  More precisely, the normalization
defect of the lifted image determines the resolution defect of the
Stein factor, and two exceptional cycles give a further inequality.
All these assertions require only ampleness of the curve.

Let $k$ be algebraically closed of characteristic two, let $X$ be a
smooth integral projective surface, and let $C\subset X$ be a smooth
integral ample divisor.  Put $d=C^2>0$.  Let $M\ne\OO_X$ be
algebraically trivial of exact order two, with $M|_C\simeq\OO_C$.
Suppose its logarithmic form has exact order two along $C$ and
\[
 2g(C)-2=3d.
\]
Assume that $A=\operatorname{Alb}(X)$ has dimension two and that
the Albanese morphism $\alpha:X\to A$ is generically finite of
degree two.  Proposition~\ref{prop:smallest-endpoint-albanese-degree-two}
gives
\begin{equation}\label{eq:strict-canonical-prior-numerics}
 \begin{gathered}
 d=4\chi(\OO_X),\qquad K_X\cdot C=2d,\qquad
 K_X\ \text{nef and big},\\
 2d\le K_X^2<3d,\qquad c_2(X)=3d-K_X^2>0.
 \end{gathered}
\end{equation}
Let $\pi:V\to X$ be the normalized torsor and let
$\beta:V\to A_V=\operatorname{Alb}(V)$ be its proper
birational Albanese morphism.
The global conductor and Albanese quotient results give
\begin{equation}\label{eq:strict-canonical-cover-data}
 \omega_V\simeq\pi^*(\omega_X(-C)),\qquad
 \alpha\pi=F_{A_V/k}\beta,\qquad A_V^{(2)}\simeq A.
\end{equation}
Here $V$ is normal and Gorenstein, and $\pi$ is finite flat of
degree two.

Retain the image notation of
Proposition~\ref{prop:smallest-endpoint-albanese-surface-defect}.
Thus $D_A=\alpha(C)$ has its reduced integral structure, $B_0$ is
its normalization, and $b=g(B_0)$.
Write $M=\alpha^*P$ and let $\rho:A'\to A$ be the degree-two
$\mu_2$-isogeny defined by $P$.
The map $B_0\to A$ lifts to $A'$; let $D'$ be its integral image and put
\begin{equation}\label{eq:strict-canonical-image-defect-definition}
 \epsilon=p_a(D')-b\ge0.
\end{equation}
This is the defect of the image in $A'$, not of an image in $A_V$.
No upper bound on $\epsilon$ is imposed.

\begin{theorem}[The canonical interval, Stein defect and image divisor]
\label{thm:smallest-endpoint-strict-canonical}
Under these hypotheses,
\begin{equation}\label{eq:strict-canonical-interval}
 2d<K_X^2<3d,\qquad 1\le c_2(X)\le d-1.
\end{equation}
Write the Stein factorization of $\alpha$ as
\[
 X\xrightarrow{\,r\,}X_0\xrightarrow{\,a_0\,}A.
\]
Then $X_0$ is normal and Gorenstein, $a_0$ is a finite flat
double cover, and $r$ is proper birational.
The quotient $Q=a_{0*}\OO_{X_0}/\OO_A$ is invertible.
Set
\[
 \mathcal L=Q^{-1},\qquad
 t=\operatorname{length}(R^1r_*\OO_X),\qquad
 D=K_X-r^*K_{X_0},
\]
using compatible canonical divisors.
The sheaf defining $t$ has finite support.
There is a nonzero effective exceptional divisor $Z$ such that
\begin{equation}\label{eq:strict-canonical-discrepancy}
 D=-Z,\qquad D^2<0,\qquad p_a(Z)=1.
\end{equation}

The map $C\to B_0$ is an isomorphism, and
\begin{equation}\label{eq:strict-canonical-image-class}
 D_A\equiv2\mathcal L
\end{equation}
numerically on $A$.
Let $E_X$ be the norm divisor constructed in the proof; it is
effective, has full exceptional support, and satisfies
$2K_X\sim2C+E_X$.
Then the actual divisor equality is
\begin{equation}\label{eq:strict-canonical-image-pullback}
 \alpha^*D_A=2C+E_X+2Z.
\end{equation}
The two defects satisfy
\begin{equation}\label{eq:strict-canonical-stein-defect}
 K_X^2=d+4t+Z^2,\qquad 2t=d+\epsilon.
\end{equation}
In particular $\epsilon$ is even.
Moreover,
\begin{equation}\label{eq:strict-canonical-double-cover-line}
 \omega_{X_0}\simeq a_0^*\mathcal L,\qquad
 \mathcal L^2=2\bigl(\chi(\OO_X)+t\bigr)
             =\frac{3d}{2}+\epsilon,
\end{equation}
and
\begin{equation}\label{eq:strict-canonical-image-cycle}
 C\cdot Z=d+2\epsilon,\qquad
 Z^2=-c_2(X)-2\epsilon.
\end{equation}

The exceptional intersection form gives
\begin{equation}\label{eq:strict-canonical-defect-tradeoff}
 4c_2(X)t\ge d^2,\qquad
 t\ge\max\left\{\frac d2,
          \left\lceil\frac{d^2}{4c_2(X)}\right\rceil\right\}.
\end{equation}
Equivalently, $2c_2(X)(d+\epsilon)\ge d^2$.
Equality in $4c_2(X)t\ge d^2$ holds precisely when
\begin{equation}\label{eq:strict-canonical-tradeoff-equality}
 (d-c_2(X))E_X=2c_2(X)Z
\end{equation}
as real exceptional divisors.
The conditions $\epsilon=0$, $t=d/2$, and smoothness of $D'$
are equivalent; in that case $c_2(X)\ge d/2$.
In particular $X_0$ has a nonrational singularity and $\alpha$
is not finite.
\end{theorem}

The canonical formula for a normal finite flat double cover is
recalled in \cite[Section~2.2, equations~(2.9)--(2.15),
pp.~8--9]{ContiPositiveSeveriComparison}.
The algebra calculation below makes its nonsplit inseparable case
explicit. The numerical range here lies above
$8\chi(\OO_X)$; the characteristic-two Severi comparison in
\cite[Theorem~8.4, p.~51]{ContiPositiveSeveriComparison} concerns
$K_X^2<(9/2)\chi(\OO_X)$ with an additional condition.

\begin{proof}
\par\smallskip\noindent\textbf{The full exceptional support of the norm divisor.}
Pull back a nowhere-zero two-form on $A_V$ by $\beta$.  On the
regular locus of $V$ this is a nonzero regular canonical section;
reflexivity of $\omega_V$ extends it over the omitted
codimension-two set.  Let $E_V$ be its effective zero divisor.

Its support is exactly the set of $\beta$-exceptional curves.
A proper birational map to a normal surface is an isomorphism
in codimension one away from its contracted curves, so there is
no zero elsewhere.  Conversely, at the generic point of a
contracted curve $F$, the surface $V$ is regular.  Two local
parameters at $\beta(F)$ restrict to constants on $F$, so their
pulled-back differentials lie in its conormal line.  Their wedge
vanishes.  Thus every contracted curve occurs in $E_V$ with
positive coefficient.

The finite-flat norm of this section, using
\eqref{eq:strict-canonical-cover-data}, gives
\begin{equation}\label{eq:strict-canonical-full-support-norm}
 2K_X\sim2C+E_X,\qquad E_X\ge0.
\end{equation}
The support of $E_X$ is the image of the support of $E_V$.
Indeed, the determinant of multiplication by a section is a unit
exactly when that section is a unit in the finite locally free
algebra.  This observation does not require a separable norm.
For a curve $F\subset V$, finiteness of $\pi$ makes
$E=\pi(F)$ a curve.  The Frobenius relation and finiteness of
Frobenius give
\[
 \beta(F)\text{ is a point}
 \quad\Longleftrightarrow\quad
 \alpha(E)\text{ is a point}.
\]
Every $\alpha$-exceptional curve has a curve above it under $\pi$.
Consequently $E_X$ contains every $\alpha$-exceptional curve,
each with positive coefficient, and has no other component.

\par\smallskip\noindent\textbf{The exceptional intersection form.}
Choose an ample divisor $H$ on $A$ and put $P=\alpha^*H$.
Then
\[
 P^2=2H^2>0,\qquad P\cdot E=0
\]
for every exceptional curve $E$.  Hodge index gives negative
square for every nonzero effective exceptional divisor.
It also gives negative definiteness for signed exceptional
divisors, as follows.  If a nonzero such divisor $U$ were
numerically zero, write $U=U_+-U_-$ with disjoint nonnegative
coefficient supports.  When $U_+\ne0$, one would have
\[
 U_+^2=U_+\cdot U_-\ge0,
\]
contrary to Hodge index applied to $U_+$.
When $U_+=0$, the divisor $-U$ is effective, nonzero and
numerically trivial, contrary to its positive intersection
with an ample divisor on $X$.  Thus a nonzero signed
exceptional divisor is never numerically zero.  Hodge index
on the orthogonal complement of $P$ proves negative definiteness.

\par\smallskip\noindent\textbf{The finite double cover.}
Stein factorization gives $r_*\OO_X=\OO_{X_0}$ and a finite
map $a_0:X_0\to A$
\cite[Theorem~37.53.4, Tag~03GX]{Stacks}.
Since $X$ is normal, $X_0$ is the normalization of $A$ in $k(X)$.
Thus $X_0$ is normal, $a_0$ has generic degree two, and $r$ is
birational.  Normal surfaces are Cohen--Macaulay.  Miracle
flatness therefore makes $a_0$ flat over the regular surface
$A$ \cite[Lemma~10.128.1, Tag~00R4]{Stacks}.
Put $\mathcal B=a_{0*}\OO_{X_0}$, a locally free algebra of rank two.

The constants form a line subbundle: over a local ring of $A$,
the element $1$ is nonzero modulo the maximal ideal and extends
to a basis of $\mathcal B$.  Hence
\begin{equation}\label{eq:strict-canonical-cover-algebra}
 0\longrightarrow\OO_A\longrightarrow\mathcal B
 \longrightarrow Q\longrightarrow0
\end{equation}
has invertible quotient.  We do not need this sequence to split.
Set $\mathcal L=Q^{-1}$.

For the canonical formula, choose a local basis $1,z$ with
$z^2=az+b$ and let $\lambda:\mathcal B\to\OO_A$ take the
coefficient of $z$.  The pairing $(u,v)\mapsto\lambda(uv)$
has matrix
\[
 \begin{pmatrix}0&1\\1&a\end{pmatrix},
\]
whose determinant is a unit in every characteristic.
Under a change $z_i=c_{ij}z_j+b_{ij}$, its functional changes
by $\lambda_i=c_{ij}^{-1}\lambda_j$.  The pairings therefore
glue to an isomorphism of $\mathcal B$-modules
\[
 \mathcal Hom_{\OO_A}(\mathcal B,\OO_A)
 \simeq\mathcal B\otimes Q^{-1}.
\]
Finite duality identifies the left side with the pushforward
of the relative dualizing sheaf
\cite[Sections~49.5 and~49.15, Tags~0FKW and~0DWM]{Stacks}.
Since $\omega_A\simeq\OO_A$, this proves that $X_0$ is
Gorenstein and
\begin{equation}\label{eq:strict-canonical-finite-identities}
 \omega_{X_0}\simeq a_0^*\mathcal L,\qquad
 \chi(\OO_{X_0})=\frac{\mathcal L^2}{2},\qquad
 K_{X_0}^2=2\mathcal L^2=4\chi(\OO_{X_0}).
\end{equation}
The Euler identity follows from
\eqref{eq:strict-canonical-cover-algebra} and Riemann--Roch on
the abelian surface, where $K_A=0$ and $\chi(\OO_A)=0$.
The intersection identity follows from finite projection.
Thus these formulas apply also to a nonsplit inseparable
double-cover algebra.

\par\smallskip\noindent\textbf{Cohomology of exceptional $(-2)$-cycles.}
We need the following vanishing: if every exceptional component
of $r$ is a smooth rational $(-2)$-curve, then
\begin{equation}\label{eq:strict-canonical-cycle-vanishing}
 R^1r_*\OO_X=0.
\end{equation}
To prove it, let $G>0$ be any effective exceptional cycle.
Negative definiteness gives $G^2<0$, so some component $E$
appearing in $G$ satisfies $G\cdot E<0$.
The Cartier-divisor sequence
\[
 0\longrightarrow\OO_E(-(G-E))
 \longrightarrow\OO_G\longrightarrow\OO_{G-E}
 \longrightarrow0
\]
is exact, including when $E$ has multiplicity in $G$.
The degree of its first term is
\[
 -(G-E)\cdot E=-G\cdot E-2\ge-1.
\]
Its first cohomology vanishes on $E\simeq\PP^1$.
Induction on the sum of the coefficients of $G$ proves
$H^1(G,\OO_G)=0$ for every effective exceptional cycle.

The map $r$ is an isomorphism away from finitely many points
and has connected fibres.  A nontrivial fibre has support
equal to the union of its exceptional curves: connectedness
excludes an isolated point component outside that union.
Over the local ring at such a point, let $I$ be the extended
maximal ideal and let $J=\OO_X(-G_{\mathrm{red}})$, where
$G_{\mathrm{red}}$ is the reduced exceptional fibre.
These ideals have the same radical, so their powers are cofinal
by Noetherianity and quasi-compactness.
For any $n$, choose $m,N$ with
\[
 I^N\subset J^m\subset I^n.
\]
Then the map
\[
 H^1(\OO_X/I^N)\longrightarrow H^1(\OO_X/I^n)
\]
factors through
$H^1(\OO_X/J^m)=H^1(\OO_{mG_{\mathrm{red}}})=0$.
Its inverse limit is therefore zero.
Formal functions identifies that limit with the completed
stalk of $R^1r_*\OO_X$
\cite[Theorem~30.20.5, Tag~02OC]{Stacks}.
The stalk is a finite module supported at the closed point;
its zero completion makes it zero.  This proves
\eqref{eq:strict-canonical-cycle-vanishing}.
No normal-crossing hypothesis on the exceptional divisor was used.

\par\smallskip\noindent\textbf{Strictness of the canonical bound.}
Suppose that $K_X^2=2d$.  Intersecting
\eqref{eq:strict-canonical-full-support-norm} with $K_X$ gives
\[
 K_X\cdot E_X=2K_X^2-2K_X\cdot C=0.
\]
By nefness and the full exceptional support of $E_X$,
$K_X\cdot E=0$ for every $\alpha$-exceptional curve.
These are exactly the exceptional curves of $r$, since
$a_0$ is finite.  Each has negative square, and adjunction gives
\[
 2p_a(E)-2=K_X\cdot E+E^2=E^2<0.
\]
Thus $p_a(E)=0$ and $E^2=-2$.
The normalization genus and normalization defect of an integral
proper curve are nonnegative and sum to its arithmetic genus.
Hence each $E$ is smooth rational, and
\eqref{eq:strict-canonical-cycle-vanishing} applies.

The divisor $K_{X_0}$ is Cartier by
\eqref{eq:strict-canonical-finite-identities}.
For compatible canonical divisors,
$D=K_X-r^*K_{X_0}$ is exceptional, and
$D\cdot E=K_X\cdot E=0$ for every exceptional $E$.
Negative definiteness forces $D=0$.
Thus $K_X^2=K_{X_0}^2$.
Normality gives $r_*\OO_X=\OO_{X_0}$, and higher direct
images above degree one vanish by fibre dimension.
Together with \eqref{eq:strict-canonical-cycle-vanishing},
this gives $\chi(\OO_X)=\chi(\OO_{X_0})$.
Equation~\eqref{eq:strict-canonical-finite-identities} now yields
\[
 K_X^2=4\chi(\OO_X)=d,
\]
contrary to $K_X^2=2d$ and $d>0$.
The prior interval \eqref{eq:strict-canonical-prior-numerics}
therefore gives \eqref{eq:strict-canonical-interval}.
In particular, some exceptional curve has positive intersection
with $K_X$; otherwise the same argument would give $K_X^2=d$.

\par\smallskip\noindent\textbf{The discrepancy and singularity defect.}
Retain $D=K_X-r^*K_{X_0}$ without the equality assumption.
It is exceptional and satisfies
\[
 D\cdot E=K_X\cdot E\ge0
\]
for every exceptional curve.
If its positive part $D_+$ were nonzero, write
$D=D_+-D_-$ with disjoint supports.  Then
\[
 0\le D_+\cdot D=D_+^2-D_+\cdot D_-<0,
\]
using negative definiteness and the nonnegative intersection
of distinct effective curves.  This contradiction proves $D\le0$.
If $D=0$, every exceptional curve again has zero canonical
intersection.  The preceding cycle argument would give
$R^1r_*\OO_X=0$ and $K_X^2=d$, contradicting
\eqref{eq:strict-canonical-prior-numerics}.
Thus $D=-Z\ne0$, with $Z$ effective and exceptional, and $D^2<0$.

The sheaf $R^1r_*\OO_X$ has finite support.  Leray, with the
degree-one direct image contributing with negative sign, gives
\[
 \chi(\OO_X)=\chi(\OO_{X_0})-t.
\]
The Cartier pullback of $K_{X_0}$ is orthogonal to $D$.
Using \eqref{eq:strict-canonical-finite-identities} and
$4\chi(\OO_X)=d$, we obtain
\[
 \begin{split}
 K_X^2
 &=K_{X_0}^2+D^2\\
 &=4\chi(\OO_X)+4t+D^2
 =d+4t+D^2.
 \end{split}
\]
Since $D^2<0$ and $K_X^2>2d$, one has
$t>(K_X^2-d)/4>0$.
Thus $R^1r_*\OO_X\ne0$, so $X_0$ has a nonrational
singularity and $r$ is not an isomorphism.
Consequently $\alpha$ is not finite.
The same Euler calculation gives
\[
 \mathcal L^2
 =2\chi(\OO_{X_0})
 =2\bigl(\chi(\OO_X)+t\bigr).
\]
Since $Z=-D$ is exceptional, $K_X\cdot Z=-Z^2$, and adjunction
for the effective Cartier cycle gives
\[
 p_a(Z)=1+\frac{Z^2+K_X\cdot Z}{2}=1.
\]

\par\smallskip\noindent\textbf{The degree of the distinguished curve.}
By Proposition~\ref{prop:smallest-endpoint-albanese-quotient},
the Albanese differential has generic rank one.
The extension $k(X)/k(A)$ has degree two and is not separable,
so it is purely inseparable of exponent one.
On a normal affine ring $R$ of $A$, let $S$ be its integral
closure in $k(X)$.
Every $s\in S$ has $s^2\in k(A)$ integral over $R$, hence
$s^2\in R$.
There is therefore a unique prime of $S$ over each prime of $R$:
if two primes have the same contraction, membership of $s$ in
one puts $s^2$ in that contraction and then puts $s$ in the other.
The residue-field extensions are purely inseparable as well.

By Proposition~\ref{prop:smallest-endpoint-albanese-forms},
the map $C\to B_0$ is finite and separable; in particular $C$
is not contracted.
The proper birational map $r$ is an isomorphism at the generic
point of $C_0=r(C)$, because $X_0$ is normal.
Thus $k(C)=k(C_0)$, and the preceding description of the finite
Stein cover makes $k(C)/k(D_A)$ purely inseparable.
It is also separable, and hence has degree one.
A finite birational map between the smooth curves $C$ and $B_0$
is an isomorphism.  In particular,
\begin{equation}\label{eq:strict-canonical-curve-degree}
 b=g(C),\qquad \alpha_*[C]=[D_A].
\end{equation}

We also need the coefficient of $C$ in $\alpha^*D_A$.
Localize $A$ at the generic point of $D_A$.
Its local ring $R$ is a discrete valuation ring, and its integral
closure $S$ in $k(X)$ is a local normal one-dimensional ring,
hence a discrete valuation ring.
The finite flat Stein cover makes $S$ free of rank two over $R$.
Its residue-field degree is one, by
\eqref{eq:strict-canonical-curve-degree}.
If a uniformizer of $R$ has valuation $e$ in $S$, the filtration
by powers of a uniformizer of $S$ gives
\[
 2=\dim_{\kappa(R)}S/\mathfrak m_RS
   =e[\kappa(S):\kappa(R)]=e.
\]
This argument applies to the purely inseparable cover directly.
The same valuation occurs along $C$ since $r$ is an isomorphism
at its generic point.
The unique nonexceptional curve over $D_A$ is $C$, so
\begin{equation}\label{eq:strict-canonical-image-initial-pullback}
 \alpha^*D_A=2C+E_b
\end{equation}
for some effective exceptional divisor $E_b$.

\par\smallskip\noindent\textbf{The image class and exceptional correction.}
For any Cartier divisor $H$ on $A$, exceptional divisors have
zero intersection with $\alpha^*H$.
The canonical identity and proper intersection projection give
\[
 K_X\cdot\alpha^*H
 =\alpha^*\mathcal L\cdot\alpha^*H
 =2\mathcal L\cdot H.
\]
Here projection holds for a proper morphism, including one
with contracted curves
\cite[Lemmas~42.26.3--42.26.4, Tags~02ST and~02SU]{Stacks}.
Intersecting \eqref{eq:strict-canonical-full-support-norm}
with $\alpha^*H$ and using
\eqref{eq:strict-canonical-curve-degree} therefore gives
\[
 4\mathcal L\cdot H
 =2K_X\cdot\alpha^*H
 =2C\cdot\alpha^*H
 =2D_A\cdot H.
\]
Every integral curve on the smooth surface $A$ is Cartier.
Thus $D_A\equiv2\mathcal L$, proving
\eqref{eq:strict-canonical-image-class}.

Now \eqref{eq:strict-canonical-image-initial-pullback}, the norm
relation, and $K_X\equiv\alpha^*\mathcal L-Z$ give
\[
 E_b-E_X-2Z\equiv0.
\]
This is a signed exceptional divisor.
Negative definiteness forces it to vanish as a divisor.
Consequently $E_b=E_X+2Z$, proving
\eqref{eq:strict-canonical-image-pullback}.
Only numerical equivalence on $A$ has been used; the pullback
formula on $X$ is an equality of divisors.

\par\smallskip\noindent\textbf{The Stein and image defects.}
The image formula in
Proposition~\ref{prop:smallest-endpoint-albanese-surface-defect},
together with $b=g(C)$, yields
\[
 D_A^2=4(b-1+\epsilon)=6d+4\epsilon.
\]
The numerical image class and the finite-cover Euler identity give
\[
 D_A^2=4\mathcal L^2
       =8\bigl(\chi(\OO_X)+t\bigr)=2d+8t.
\]
Comparing them proves $2t=d+\epsilon$.
Since $4\mid d$, the integer $\epsilon$ is even.
Its nonnegativity gives $t\ge d/2$, and substitution gives
\[
 \mathcal L^2=\frac{3d}{2}+\epsilon.
\]
This proves the defect and line identities in the statement.

The norm relation gives $C\cdot E_X=2d$.
Intersecting the exact pullback formula with $C$ and applying
\eqref{eq:strict-canonical-curve-degree}, we obtain
\[
 6d+4\epsilon=D_A^2
 =2d+C\cdot E_X+2C\cdot Z=4d+2C\cdot Z.
\]
Thus $C\cdot Z=d+2\epsilon$.
Also
\[
 Z^2=K_X^2-d-4t
    =K_X^2-3d-2\epsilon=-c_2(X)-2\epsilon,
\]
proving \eqref{eq:strict-canonical-image-cycle}.
The equivalence of $\epsilon=0$ with smoothness of $D'$ was
proved in Proposition~\ref{prop:smallest-endpoint-albanese-surface-defect};
the defect identity makes it equivalent to $t=d/2$.

\par\smallskip\noindent\textbf{The exceptional-cycle inequality.}
Put $c=c_2(X)$ and consider the effective real exceptional divisors
\[
 U=\frac12E_X,\qquad W=Z+\frac12E_X.
\]
Their numerical classes are $U\equiv K_X-C$ and
$W\equiv\alpha^*\mathcal L-C$.
Hence
\[
 U^2=(K_X-C)^2=K_X^2-3d=-c.
\]
The image class and the curve degree give
\[
 \alpha^*\mathcal L\cdot C
 =\mathcal L\cdot D_A
 =2\mathcal L^2=(\alpha^*\mathcal L)^2.
\]
Consequently
\[
 W^2=d-2\mathcal L^2=-4t.
\]
Since $U$ is exceptional,
\[
 U\cdot W=-U\cdot C=-(K_X\cdot C-d)=-d.
\]
The negative of the intersection form is a positive-definite inner
product on the real exceptional space.
Cauchy--Schwarz gives $d^2\le4ct$.
Combining this with $t\ge d/2$ and integrality of $t$ proves
\eqref{eq:strict-canonical-defect-tradeoff}.
Replacing $2t$ by $d+\epsilon$ gives its equivalent form.

Equality holds precisely when $W=(d/c)U$, since the inner
products of $U$ with itself and with $W$ are $c$ and $d$.
This is equivalent to $(d-c)E_X=2cZ$, proving
\eqref{eq:strict-canonical-tradeoff-equality}.
Finally, when $\epsilon=0$, the inequality becomes
$2cd\ge d^2$, and hence $c\ge d/2$.
\end{proof}

\begin{remark}[Degree sixteen and scope]
\label{rem:smallest-endpoint-strict-degree-sixteen}
If $C$ is very ample and $d=16$, then
Corollary~\ref{cor:smallest-endpoint-albanese-very-ample-degree}
gives
\[
 h^0(X,\OO_X(C))=6,\qquad h^0(C,\OO_C(C))=5.
\]
Theorem~\ref{thm:smallest-endpoint-strict-canonical} requires
\[
 33\le K_X^2\le47,\qquad 1\le c_2(X)\le15,\qquad
 t\ge\max\left\{8,\left\lceil\frac{64}{c_2(X)}\right\rceil\right\},
\]
and
\[
 \epsilon=2t-16,\qquad \mathcal L^2=24+\epsilon.
\]
For example, $c_2(X)=1,2,4,8$ require respectively
$t\ge64,32,16,8$.
If $D'$ is smooth, then $t=8$ and $8\le c_2(X)\le15$.
These are necessary conditions, with no existence assertion.

All conclusions of the theorem require only ampleness of $C$;
very ampleness is used here solely for the degree and section counts.
The surface $X_0$ is the finite Stein factor of $\alpha$,
not the foliation quotient $V^{(2)}$.
The cycle $Z$ may be reducible or nonreduced; $p_a(Z)=1$
does not assert that an exceptional component is a smooth elliptic
curve.
The numerical identity $D_A\equiv2\mathcal L$ does not assert
linear equivalence.
The normalization defect $\epsilon$ is not bounded above, and
the exceptional-cycle inequality does not exclude the entire
degree-two branch.
The image-pattern examples in
Remark~\ref{rem:smallest-endpoint-albanese-image-pattern}
do not satisfy all the endpoint surface hypotheses.
\end{remark}
